%----------------------------------------------------------------------------------------
% PACKAGES AND THEMES
%----------------------------------------------------------------------------------------
\documentclass[aspectratio=169,xcolor=dvipsnames]{beamer}

\usetheme{SimplePlus}

\usepackage{hyperref}
\usepackage{helvet}
\usepackage{graphicx}
\usepackage{booktabs}
\usepackage{ragged2e}

\setbeamerfont{frametitle}{size=\Large}
\setbeamerfont{normal text}{size=\large}

%----------------------------------------------------------------------------------------
% TITLE PAGE INFORMATION
%----------------------------------------------------------------------------------------

%----------------------------------------------------------------------------------------
\begin{document}

%------------------------------------------------
% TITLE PAGE
%------------------------------------------------
\begin{frame}[plain]

\centering

%------------------------------------------------
% COLLEGE HEADER
%------------------------------------------------
\includegraphics[
    width=\paperwidth,
    height=0.19\paperheight,
    keepaspectratio
]{Established 2001 copy.png}

\vspace{0.20cm}

%------------------------------------------------
% PROJECT TITLE
%------------------------------------------------
{\fontsize{11}{13}\selectfont
\textbf{
An Event-Driven Geospatial Push Methodology for \\
Urban Food Waste Reduction 
Utilizing Dynamic Radius Computation and Real-Time Urgency Scoring
}
}

\vspace{0.28cm}

%------------------------------------------------
% STUDENT NAME
%------------------------------------------------
{\fontsize{12}{14}\selectfont
\textbf{Vijayprasath G - 311523205061}
}

\vspace{0.25cm}

%------------------------------------------------
% GUIDE
%------------------------------------------------
{\fontsize{11}{12}\selectfont
\textbf{Guided by:}\\
\textbf{Mohan Raj Vijayan}
}

\vspace{0.08cm}

%------------------------------------------------
% GUIDE QUALIFICATION
%------------------------------------------------
{\fontsize{9}{10}\selectfont
M.E (CSE), (Ph.D.), M.B.A (PM), M.Music,\\
M.Sc. (Psy), M.A (Philo), M.A (French),\\
Dip. (Arch \& Epi)
}

\vspace{0.16cm}

%------------------------------------------------
% DESIGNATION
%------------------------------------------------
{\fontsize{11}{12}\selectfont
\textbf{Assistant Professor}\\
Department of Information Technology
}

\vspace{0.08cm}

%------------------------------------------------
% COLLEGE
%------------------------------------------------
{\fontsize{9}{11}\selectfont
Meenakshi Sundararajan Engineering College, Chennai
}

\vspace{0.12cm}

%------------------------------------------------


\end{frame}

%------------------------------------------------
%----------------------------------------------------------------------------------------
% TITLE SLIDE
%----------------------------------------------------------------------------------------

%----------------------------------------------------------------------------------------
% TABLE OF CONTENTS
%----------------------------------------------------------------------------------------
\begin{frame}{Table of Contents}
\tableofcontents
\end{frame}

%----------------------------------------------------------------------------------------
\section{Problem Statement}
\begin{frame}{Problem Statement}
\justifying
Every day, local food vendors discard massive amounts of unsold, perfectly good inventory at closing time, leading to significant economic and environmental losses. Current food-rescue platforms attempt to mitigate this but suffer from critical limitations: they rely on static, ``pull-based'' architectures where consumers must actively remember to search for deals, and they often obscure items in opaque ``mystery bags'' that deter user participation. 

While recent frameworks like SustainNet (Krishnan et al., 2025) have proven the value of PostGIS geospatial matching, a significant technical gap remains. There is a lack of an event-driven, ``push-based'' system that actively alerts nearby users the exact moment surplus becomes available. Furthermore, existing solutions fail to provide frictionless vendor uploading (via Vision AI) and lack dynamic, time-decaying pricing models (Dutch Auctions) necessary to algorithmically incentivize immediate food rescue in the final critical minutes before closing.
\end{frame}


%----------------------------------------------------------------------------------------
\section{Social Relevance}
\begin{frame}{Social Relevance}
\justifying

\begin{itemize}
    \item \textbf{Reduces Food Waste:} Actively connects unsold food with nearby consumers in the final critical minutes before closing, preventing meals from being discarded.

    \item \textbf{Supports Local Vendors:} Enables small businesses to recover ingredient costs from surplus inventory instead of absorbing complete losses daily.

    \item \textbf{Promotes Transparency:} Replaces opaque ``mystery bags'' with exact, transparent item listings so consumers know precisely what they are purchasing.

    \item \textbf{Privacy-First Design:} Uses location solely for real-time matching without storing any consumer location history, ensuring strict data privacy compliance.

    \item \textbf{Strengthens Community Networks:} Creates hyper-local connections between neighborhood vendors and residents, fostering sustainable urban food ecosystems.
\end{itemize}

\end{frame}
%----------------------------------------------------------------------------------------

\section{Abstract}
\begin{frame}{Abstract}
\justifying

This project proposes a hyper-local, event-driven food rescue platform designed to eliminate end-of-day food waste. The objective is to actively connect food vendors with surplus inventory to nearby consumers in real-time, moving away from passive browsing models. The system employs an Ephemeral Geospatial Push Protocol that utilizes PostGIS spatial queries, Dynamic Radius Computation, and a Real-Time Urgency Score (RUS) engine to deliver contextual push notifications. To solve the pricing problem, a Dutch Auction Pricing ML model is integrated — a regression-based engine that predicts the time-to-claim at varying discount levels and automatically schedules price reductions as the store approaches closing, maximizing vendor revenue while guaranteeing a sale. A highly targeted system that accelerates food claim rates while ensuring strict data ephemerality and preventing regular sales cannibalization. This event-driven approach provides a scalable, privacy-first architecture that significantly reduces urban food waste.

\end{frame}

%----------------------------------------------------------------------------------------
\section{Objectives and Sustainability Alignment}

\begin{frame}{Objectives and Sustainability Alignment}
\justifying

\begin{columns}[t]

%------------------------------------------------
% Objectives
%------------------------------------------------
\column{0.48\textwidth}

\textbf{Objectives}

\begin{itemize}
    \item To develop an event-driven geospatial push framework for real-time food rescue matching.

    \item To integrate PostGIS spatial queries, dynamic radius computation, and Real-Time Urgency Scoring for accurate contextual delivery.

    \item To build an ephemeral data architecture that ensures all deal data vanishes once the vendor closes.
\end{itemize}

%------------------------------------------------
% Sustainability Alignment
%------------------------------------------------
\column{0.48\textwidth}

\textbf{Sustainability Alignment -- SDG 12}

\begin{itemize}
    \item \textbf{12.3: Responsible Consumption and Production} -- Reduce food waste at the retail level through active surplus redistribution.

    \item \textbf{12.5: Substantially Reduce Waste Generation} -- Prevent end-of-day food disposal through timely vendor-consumer matching.

    \item \textbf{12.8: Information and Awareness} -- Promote transparent, informed food purchasing over opaque mystery-bag models.
\end{itemize}

\end{columns}

\end{frame}

%----------------------------------------------------------------------------------------
\section{Scope}
\begin{frame}{Scope}
\justifying

\begin{itemize}
    \item Support local food vendors (bakeries, restaurants, cafes) by providing a PWA to instantly post exact, transparent inventory right before closing time.

    \item Actively push deals to consumers within a dynamically calculated geographic radius using PostGIS and Firebase Cloud Messaging (FCM).

    \item Implement Dynamic Price Degradation (Dutch Auction) to maximize vendor revenue and eliminate the cannibalization of regular daytime sales.

    \item Ensure all deal data is strictly ephemeral and guarantee true zero-waste via an automated charity cascade (NGO Handoff) for unsold items 15 minutes before closing.

    \item Provide vendor analytics dashboards showing rescue rates, revenue recovered, and optimal posting times.
\end{itemize}

\end{frame}
%----------------------------------------------------------------------------------------
\section{Novelty and Innovation Statement}

\begin{frame}{Novelty and Innovation Statement}
\justifying

\begin{columns}[t]

%------------------------------------------------
% Novelty
%------------------------------------------------
\column{0.48\textwidth}

\textbf{Novelty}

\begin{itemize}
    \item Introduces an \textbf{Ephemeral Geospatial Push Protocol} that actively notifies users instead of waiting for them to browse.

    \item Combines \textbf{PostGIS spatial queries, Dutch Auction pricing, and Real-Time Urgency Scoring} into a single event-driven pipeline.

    \item Enforces \textbf{true zero-waste} via an automated NGO Handoff cascade, and \textbf{strict data ephemerality} when the store closes.
\end{itemize}

%------------------------------------------------
% Innovation
%------------------------------------------------
\column{0.48\textwidth}

\textbf{Innovation}

\begin{itemize}
    \item Uses \textbf{ST\_DWithin on true geography (sphere)} instead of approximate bounding box calculations for spatial precision.

    \item Implements a \textbf{5-factor Real-Time Urgency Score (RUS)} combining time pressure, scarcity, claim speed, distance, and weather.

    \item Employs \textbf{anti-cannibalization algorithms} that exclude repeat claimers, turning the platform into a customer acquisition tool.
\end{itemize}

\end{columns}

\end{frame}

%----------------------------------------------------------------------------------------
%========================================================
% LITERATURE REVIEW
%========================================================

\section{Literature Review}

%--------------------------------------------------------
% PAGE 1
%--------------------------------------------------------
\begin{frame}{Literature Review (1/5)}
\scriptsize
\renewcommand{\arraystretch}{1.25}
\setlength{\tabcolsep}{2.5pt}

\begin{tabular}{|c|p{1.7cm}|p{3.6cm}|p{3.1cm}|p{3.2cm}|}
\hline
\textbf{No.} &
\textbf{Author (Year)} &
\textbf{Paper Title} &
\textbf{Mechanism Used} &
\textbf{Drawbacks} \\
\hline

1 &
Krishnan S. et al. (2025) &
SustainNet: A Cloud-Based Surplus Food Redistribution Network with Geospatial Matching &
PostGIS, React, Cloud Microservices, JWT &
No push notifications; no urgency scoring; no ephemeral data design \\
\hline

2 &
Suresh G. H. et al. (2025) &
Food Waste Management and Prediction Using Machine Learning &
Random Forest, XGBoost, k-means &
Relies on synthetic data; limited to specific food categories \\
\hline

3 &
Ahmed S. et al. (2025) &
Sustainable Farming with ML Solutions for Minimizing Food Waste &
Random Forest, Gradient Boosting, Time Series &
High IoT sensor cost; accuracy varies across food types \\
\hline

4 &
Okonkwo E. et al. (2025) &
Leveraging AI to Minimise Food Waste and Enhance Global Food Security &
CNN, LSTM, Reinforcement Learning &
Review/framework paper; lacks empirical validation \\
\hline

5 &
Dehghan A. et al. (2025) &
ML-Based Prediction of C/N Ratio in Municipal Organic Waste &
Random Forest, Ridge/Lasso, SVR &
Limited to municipal organic waste; small sample size \\
\hline

\end{tabular}
\end{frame}


%--------------------------------------------------------
% PAGE 2
%--------------------------------------------------------
\begin{frame}{Literature Review (2/5)}
\scriptsize
\renewcommand{\arraystretch}{1.35}
\setlength{\tabcolsep}{2.5pt}

\begin{tabular}{|c|p{1.7cm}|p{3.8cm}|p{3.2cm}|p{3.5cm}|}
\hline
\textbf{No.} &
\textbf{Author (Year)} &
\textbf{Paper Title} &
\textbf{Mechanism Used} &
\textbf{Drawbacks} \\
\hline

6 &
Singh I. et al. (2025) &
Interpretable Hybrid Deep Q-Learning Framework for IoT-Based Food Spoilage Prediction &
Deep Q-Learning, LSTM, GAN, SHAP/LIME &
Lab conditions only; limited to 3 food types \\
\hline

7 &
Yang H. \& Yu M. (2025) &
Surprise Clearance Models and Food Waste: When Mystery Bags Transfer Waste &
Game-theoretic modeling, Monte Carlo &
Analytical model; not validated with real transaction data \\
\hline

8 &
Bose D. et al. (2025) &
AI-Driven Food Surplus Redistribution: Predictive Matching Between Restaurants and NGOs &
Random Forest, Haversine, FastAPI &
Small pilot (single city); no push notification system \\
\hline

9 &
Gershengoren L. \& Bandi S. (2025) &
Consumer Motivations and Barriers in Digital Food Waste Reduction Platforms &
EFA, Multiple Linear Regression &
Self-reported survey data; cross-sectional design \\
\hline

10 &
Chen W. et al. (2025) &
Factors Influencing Purchase Intentions for Surplus Food Blind Boxes &
PLS-SEM, Confirmatory Factor Analysis &
Survey-based; does not measure actual purchase behavior \\
\hline

\end{tabular}
\end{frame}


%--------------------------------------------------------
% PAGE 3
%--------------------------------------------------------
\begin{frame}{Literature Review (3/5)}
\scriptsize
\renewcommand{\arraystretch}{1.25}
\setlength{\tabcolsep}{2.5pt}

\begin{tabular}{|c|p{1.7cm}|p{3.6cm}|p{3.1cm}|p{3.2cm}|}
\hline
\textbf{No.} &
\textbf{Author (Year)} &
\textbf{Paper Title} &
\textbf{Mechanism Used} &
\textbf{Drawbacks} \\
\hline

11 &
Vijayakumar K. et al. (2025) &
Food Wastage Analysis Framework Using Comprehensive IoT-Based Multi-Sensor Systems &
Random Forest, SVM, AWS IoT, MQTT &
High sensor cost; tested only in controlled environments \\
\hline

12 &
Torres M. et al. (2025) &
Real-Time Food Waste Monitoring Using Image Processing and IoT in Commercial Kitchens &
ResNet-50 CNN, IoT Load Cells, Node-RED &
Tested in 5 similar kitchens; privacy concerns with cameras \\
\hline

13 &
Hartono B. et al. (2025) &
Food Waste Management Strategy with Green and Digital Technology: Computer Vision &
YOLOv8, CNN, Sensor Fusion &
Requires GPU for real-time; limited generalization \\
\hline

14 &
Petersen N. et al. (2025) &
Multimodal Dataset and Framework for Non-Invasive Freshness Assessment of Perishable Produce &
Multimodal Deep Learning, CNN, Ensemble &
Small dataset (5 produce types); subjective labels \\
\hline

15 &
Rossi F. et al. (2026) &
Multichannel Gas Sensor System with TinyML for Real-Time Perishable Food Shelf-Life Prediction &
TinyML, Random Forest, ESP32 Edge &
Limited to 8 produce types; controlled gas environment only \\
\hline

\end{tabular}
\end{frame}


%--------------------------------------------------------
% PAGE 4
%--------------------------------------------------------
\begin{frame}{Literature Review (4/5)}
\scriptsize
\renewcommand{\arraystretch}{1.35}
\setlength{\tabcolsep}{3pt}

\begin{tabular}{|c|p{2.0cm}|p{3.0cm}|p{2.8cm}|p{3.0cm}|}
\hline
\textbf{No.} &
\textbf{Author (Year)} &
\textbf{Paper Title} &
\textbf{Mechanism Used} &
\textbf{Drawbacks} \\
\hline

16 &
Garg A. et al. (2025) &
Dynamic Delivery Time Prediction Using Real-Time Contextual Features &
LightGBM, XGBoost, Haversine, Real-Time API &
Single Indian metro; 15-min traffic granularity \\
\hline

17 &
Li J. et al. (2025) &
Knowledge-Enhanced Spatial-Temporal Routing Optimization for Last-Mile Food Delivery &
ST-GNN, Knowledge Graphs, Dijkstra's &
Single city; static road network graph \\
\hline

18 &
Al-Mansouri F. et al. (2026) &
Composite GiST Indexing for Spatiotemporal Query Optimization in High-Volume Logistics &
PostGIS GiST/SP-GiST, BRIN, ST\_DWithin &
Single PostgreSQL instance; gains diminish for small datasets \\
\hline

19 &
Kavoosi A. et al. (2026) &
Dynamic Pricing and Inventory Control of Perishable Products Using Deep RL &
PPO, TD3, MDP, Neural Networks &
Single retail chain; 3 product categories only \\
\hline

20 &
Muller T. et al. (2025) &
Blockchain-Based Food Traceability Framework with GDPR-Compliant Off-Chain Storage &
Hyperledger Fabric, IPFS, Zero-Knowledge Proofs &
Small pilot (50 participants); no cross-blockchain interop \\
\hline

\end{tabular}
\end{frame}


%--------------------------------------------------------
% PAGE 5
%--------------------------------------------------------
\begin{frame}{Literature Review (5/5)}
\scriptsize
\renewcommand{\arraystretch}{1.35}
\setlength{\tabcolsep}{3pt}

\begin{tabular}{|c|p{2.0cm}|p{3.0cm}|p{2.8cm}|p{3.0cm}|}
\hline
\textbf{No.} &
\textbf{Author (Year)} &
\textbf{Paper Title} &
\textbf{Mechanism Used} &
\textbf{Drawbacks} \\
\hline

21 &
Patil V. et al. (2025) &
Smart Contract-Driven Food Donation Platforms with QR-Code Verification &
Ethereum/Solidity, QR Code, ERC-20 Tokens &
High gas fees; limited pilot (15 donors, 5 NGOs) \\
\hline

22 &
Zhang H. et al. (2025) &
Explainable AI-Enhanced QR-Based Food Supply Chain Traceability &
CNN, SHAP (XAI), Hyperledger, QR APIs &
Single product category (dairy); 200-respondent survey \\
\hline

23 &
Sereda D. \& Kowalski M. (2025) &
Comparative Performance Analysis of SSE vs WebSockets in High-Concurrency Systems &
SSE, WebSocket, k6/Artillery, Prometheus &
Single-server setups; SSE is unidirectional only \\
\hline

24 &
Rahman A. et al. (2025) &
Serverless Transaction Management for Event-Driven Food Delivery Using Saga Patterns &
Saga Pattern, Apache Kafka, AWS Lambda, CQRS &
Simulated workload only; Saga compensation manually coded \\
\hline

25 &
Park J. et al. (2025) &
BE-YOLOv8: Improved YOLOv8 with Attention Mechanisms for Real-Time Food Waste Classification &
YOLOv8, CBAM Attention, LMix, ONNX &
Requires GPU for real-time; degrades on occluded items \\
\hline

\end{tabular}
\end{frame}


%----------------------------------------------------------------------------------------

%------------------------------------------------
\section{Design Methodology}

%================================================
% PAGE 1
%================================================
\begin{frame}{Design Methodology}
\justifying
\scriptsize
\setlength{\itemsep}{1pt}
\setlength{\parskip}{0pt}

\textbf{Step 1: Vendor Deal Posting}
\begin{itemize}
    \item Input: Exact inventory items, dietary tags, closing time, rescue pricing.
    \item Vendor submits via React PWA; backend validates and creates deal record.
    \item \textbf{Technique:} Zod Schema Validation + JWT Authentication.
\end{itemize}

\vspace{0.05cm}

\textbf{Step 2: Dynamic Radius Computation}
\begin{itemize}
    \item Calculate effective push radius based on city density, time urgency, item quantity, and historical claim rate.
    \item Hard limits enforced: minimum 500m, maximum 5000m.
    \item \textbf{Technique:} Multi-Factor Radius Algorithm.
\end{itemize}

\vspace{0.05cm}

\textbf{Step 3: Geospatial User Matching}
\begin{itemize}
    \item Query users physically within the computed radius using PostGIS.
    \item Filter by location freshness ($<$15 min), notification consent, and anti-cannibalization rules.
    \item \textbf{Technique:} PostGIS ST\_DWithin + Geography Indexing (GIST).
\end{itemize}

\vspace{0.05cm}

\textbf{Step 4: Urgency Scoring \& Price Degradation}
\begin{itemize}
    \item Calculate a composite 0--10 urgency score (RUS) using 5 weighted factors.
    \item Dynamically drop the price of the deal (Dutch Auction) as closing time approaches to force immediate claims.
    \item Map to labels: GOING\_NOW ($\geq$8), ACT\_SOON ($\geq$5), COMFORTABLE ($<$5).
    \item \textbf{Technique:} Multi-Factor Scoring \& Algorithmic Price Degradation.
\end{itemize}

\end{frame}

%================================================
% PAGE 2
%================================================
\begin{frame}{Design Methodology}
\justifying

\textbf{Step 5: Active Push Delivery}
\begin{itemize}
    \item Dispatch deal notifications to matched users via Server-Sent Events (SSE) for in-app delivery.
    \item Simultaneously send Firebase Cloud Messaging (FCM) for background/locked-screen delivery.
    \item \textbf{Technique:} Dual-Channel Push (SSE + FCM).
\end{itemize}

\vspace{0.15cm}

\textbf{Step 6: Claim Processing and QR Verification}
\begin{itemize}
    \item Use optimistic concurrency (atomic SQL) to prevent overselling.
    \item Generate time-limited QR code; vendor scans at pickup to confirm.
    \item Auto-release uncollected claims after 20 minutes; re-push to new users.
    \item \textbf{Technique:} Optimistic Locking + BullMQ Delayed Jobs.
\end{itemize}

\vspace{0.15cm}

\textbf{Step 7: NGO Handoff \& Ephemeral Data Sweep}
\begin{itemize}
    \item Unsold items 15 mins before closing trigger an automated webhook cascade to partner NGOs (Charity Handoff).
    \item Background workers automatically expire deals at closing time, firing SSE ``deal:expired'' events.
    \item Purge deal data 24h post-expiry; claims purged after 30 days.
    \item \textbf{Output:} Zero-Waste Guarantee \& Strict Ephemeral State Consistency.
\end{itemize}

\vspace{0.2cm}

\textbf{Overall Pipeline:}
Deal Posting $\rightarrow$ Radius Computation $\rightarrow$
Geospatial Matching $\rightarrow$ Urgency Scoring \& Price Drop $\rightarrow$
Active Push $\rightarrow$ Claim \& QR Pickup $\rightarrow$ NGO Handoff $\rightarrow$ Ephemeral Sweep

\end{frame}

%----------------------------------------------------------------------------------------
\section{Machine Learning Integration}
\begin{frame}{Machine Learning Integration}
\justifying
\scriptsize

\textbf{1. Dynamic Price \& Demand Suggester (Dutch Auction Engine)}
\begin{itemize}
    \item \textbf{Objective:} Predict the optimal Start Price, Minimum Reserve Price, and Degradation Curve.
    \item \textbf{Model:} Scikit-learn based Regression Model (Random Forest).
    \item \textbf{Input Features:} Time until close, item category, original price, weather factor, city density, and historical claim lag.
    \item \textbf{Output:} Predicted time-to-claim for varying price points.
\end{itemize}

\vspace{0.1cm}

\textbf{2. Snap-to-Post Vision AI (Food Image Recognition)}
\begin{itemize}
    \item \textbf{Objective:} Automatically identify food items from a vendor's camera photo and auto-fill the deal posting form.
    \item \textbf{Model:} Pre-trained Vision Transformer (ViT) fine-tuned on Indian food datasets via Hugging Face (\texttt{rajistics/finetuned-indian-food}, 93\%+ accuracy).
    \item \textbf{Pipeline:} Vendor snaps photo $\rightarrow$ Cloudinary CDN upload $\rightarrow$ ViT classification $\rightarrow$ Taxonomy mapper infers dietary tags \& benchmark price $\rightarrow$ Auto-triggers Dutch Auction pricing.
    \item \textbf{Fallback:} If confidence $<$ 0.50 or inference exceeds 2s timeout, form defaults to manual entry with image attachment.
\end{itemize}

\vspace{0.1cm}
\textbf{Architecture:} Both models run in a single Python FastAPI microservice, pre-warmed on startup, communicating with the Node.js backend.
\end{frame}

%----------------------------------------------------------------------------------------
\section{System Architecture}
\begin{frame}{System Architecture}
\justifying
\centering

\includegraphics[
    width=\textwidth,
    height=0.9\textheight,
    keepaspectratio
]{vijay architecture.png}

\end{frame}


%----------------------------------------------------------------------------------------
\section{Software and Hardware Tools Identified}

\begin{frame}{Software and Hardware Tools Identified}

\begin{columns}

% Left Column
\column{0.48\textwidth}
\textbf{Software Tools}
\begin{itemize}
    \item React.js (PWA)
    \item Node.js 20 LTS + Express.js
    \item PostgreSQL 16 + PostGIS 3.4
    \item Redis 7 + BullMQ
    \item Python 3.11 + FastAPI (AI Service)
    \item Hugging Face Transformers (ViT)
    \item Firebase Admin SDK (FCM)
    \item Leaflet.js + OpenStreetMap
    \item Cloudinary (Image CDN)
    \item Docker + Nginx
\end{itemize}

% Right Column
\column{0.48\textwidth}
\textbf{Hardware Requirements}
\begin{itemize}
    \item Intel Core i5/i7 Processor
    \item 8 GB RAM (Minimum)
    \item 16 GB RAM (Recommended)
    \item 256 GB SSD Storage
    \item Internet Connectivity
    \item GPS-Enabled Mobile Device
    \item Windows/Linux Operating System
\end{itemize}

\end{columns}

\end{frame}

%----------------------------------------------------------------------------------------
\section{Feasibility Analysis}
\begin{frame}{Feasibility Analysis}
\justifying
\begin{itemize}
    \item The proposed system utilizes mature, production-ready technologies such as PostgreSQL + PostGIS, Node.js, React.js, Redis, and Firebase, ensuring reliable and stable implementation.
    
    \item PostGIS provides industry-standard geospatial indexing (GIST) and precise geography calculations via ST\_DWithin, enabling accurate and efficient spatial queries at scale.
    
    \item Server-Sent Events (SSE) leverage built-in browser reconnection support and require no external WebSocket infrastructure, reducing deployment complexity.
    
    \item BullMQ with Redis provides robust delayed job scheduling for TTL sweeps, claim timeouts, and RUS recalculation, all with built-in retry and backoff mechanisms.
    
    \item The PWA architecture eliminates app store deployment barriers, enabling rapid vendor onboarding with a lightweight, installable web application accessible on any device.
\end{itemize}
\end{frame}

%----------------------------------------------------------------------------------------
\section{Innovation and Concept Validation}
\begin{frame}{Innovation and Concept Validation}
\justifying
\textbf{Innovation}
\begin{itemize}
    \item Proposes a novel Ephemeral Geospatial Push Protocol that shifts food rescue from passive browsing to active, location-aware notification delivery.
    
    \item Integrates a multi-factor Dynamic Radius Engine with PostGIS spatial queries, combined with a Dutch Auction Pricing Model to maximize vendor revenue.
    
    \item Employs a 5-factor Real-Time Urgency Score (RUS) that contextualizes each deal's criticality, and an Automated NGO Handoff to guarantee true zero-waste.
\end{itemize}

\textbf{Concept Validation}
\begin{itemize}
    \item Anti-cannibalization logic ensures the platform serves as a customer acquisition channel rather than a discount tool for existing regulars.
    
    \item Optimistic concurrency with atomic SQL and BullMQ auto-release ensures zero overselling and automatic slot recovery from no-shows.
\end{itemize}
\end{frame}

%----------------------------------------------------------------------------------------
\section{Conclusion}
\begin{frame}{Conclusion}
\justifying
\begin{itemize}

    \item Proposed an \textbf{Ephemeral Geospatial Push Protocol (EGPP)} that fundamentally shifts food rescue from passive browsing to active, location-aware push delivery.

    \item Integrates \textbf{PostGIS spatial queries, dynamic radius computation, Real-Time Urgency Scoring, SSE, and FCM} into a unified event-driven pipeline.

    \item Combines \textbf{Dutch Auction Pricing} with an \textbf{Automated NGO Handoff} to maximize vendor revenue, eliminate cannibalization, and guarantee absolute zero waste.

    \item Enforces \textbf{strict data ephemerality} where all deal data is automatically swept server-side the moment the vendor closes.

    \item Provides a \textbf{scalable, privacy-first, and vendor-friendly architecture} for sustainable hyper-local food rescue across multiple urban density tiers.

\end{itemize}
\end{frame}

%----------------------------------------------------------------------------------------
\section{Future Scope and Recommendations}
\begin{frame}{Future Scope and Recommendations}
\justifying
\textbf{Future Scope}
\begin{itemize}
    \item Integrate hyper-local delivery partners and volunteer networks for last-mile food rescue to reach consumers unable to travel.
    
    \item Expand the RUS algorithm with machine learning for predictive surplus forecasting based on historical vendor data and seasonal patterns.
    
    \item Develop API bridges to food banks and NGOs for automatic routing of un-purchased inventory after closing time.
    
    \item Extend Snap-to-Post Vision AI with multi-item batch classification and multi-language food label support for broader regional coverage.
\end{itemize}

\textbf{Recommendations}
\begin{itemize}
    \item Extend the architecture to support non-food perishable goods such as surplus local agriculture and cut flowers.
    
    \item Implement real-time social rescue features allowing neighbors to collaborate on bulk surplus pick-ups.
    
    \item Continuously update weather and density data sources for improved RUS accuracy across diverse geographies.
\end{itemize}
\end{frame}

%----------------------------------------------------------------------------------------
\section{References}

%--------------------------------------------------------
% REFERENCES PAGE 1 (Papers 1--7)
%--------------------------------------------------------
\begin{frame}{References (1/4)}
\tiny
\begin{enumerate}
    \item Krishnan S., Reddy A., Nair P., ``SustainNet: A Cloud-Based Surplus Food Redistribution Network with Geospatial Matching,'' 2025. \textbf{[BASE PAPER]}

    \item Suresh G. H., Priya R., Ramesh K., ``Food Waste Management and Prediction Using Machine Learning: A Comprehensive Analysis,'' 2025.

    \item Ahmed S., Kumar V., Zhang L., ``Sustainable Farming with Machine Learning Solutions for Minimizing Food Waste,'' 2025.

    \item Okonkwo E., Patel M., Williams J., ``Leveraging Artificial Intelligence to Minimise Food Waste and Enhance Global Food Security,'' 2025.

    \item Dehghan A., Farahani M., Rezaei H., ``Machine Learning-Based Prediction of the C/N Ratio in Municipal Organic Waste,'' 2025.

    \item Singh I., Sharma P., Gupta A., ``Interpretable Hybrid Deep Q-Learning Framework for IoT-Based Food Spoilage Prediction,'' 2025.

    \item Yang H., Yu M., ``Surprise Clearance Models and Food Waste: When Mystery Bags Transfer Waste from Store to Consumer,'' 2025.
\end{enumerate}
\end{frame}

%--------------------------------------------------------
% REFERENCES PAGE 2 (Papers 8--14)
%--------------------------------------------------------
\begin{frame}{References (2/4)}
\tiny
\begin{enumerate}
\setcounter{enumi}{7}
    \item Bose D., Chakraborty S., Das R., ``AI-Driven Food Surplus Redistribution: Predictive Matching Between Restaurants and NGOs,'' 2025.

    \item Gershengoren L., Bandi S., ``Consumer Motivations and Barriers in Digital Food Waste Reduction Platforms,'' 2025.

    \item Chen W., Liu X., Wang Y., ``Factors Influencing Purchase Intentions for Surplus Food Blind Boxes: Perceived Risk and Taste Anxiety,'' 2025.

    \item Vijayakumar K., Anand R., Priya S., ``Food Wastage Analysis Framework Using Comprehensive IoT-Based Multi-Sensor Systems,'' 2025.

    \item Torres M., Garcia L., Fernandez J., ``Real-Time Food Waste Monitoring Using Image Processing and IoT in Commercial Kitchens,'' 2025.

    \item Hartono B., Sari D., Wijaya R., ``Food Waste Management Strategy with Green and Digital Technology: Computer Vision Approaches,'' 2025.

    \item Petersen N., Andersen M., Larsen K., ``Multimodal Dataset and Framework for Non-Invasive Freshness Assessment of Perishable Produce,'' 2025.
\end{enumerate}
\end{frame}

%--------------------------------------------------------
% REFERENCES PAGE 3 (Papers 15--21)
%--------------------------------------------------------
\begin{frame}{References (3/4)}
\tiny
\begin{enumerate}
\setcounter{enumi}{14}
    \item Rossi F., Bianchi G., Moretti L., ``Multichannel Gas Sensor System with TinyML for Real-Time Perishable Food Shelf-Life Prediction,'' 2026.

    \item Garg A., Mehta S., Roy D., ``Dynamic Delivery Time Prediction Using Real-Time Contextual Features in Indian Urban Environments,'' 2025.

    \item Li J., Wang Z., Chen Y., ``Knowledge-Enhanced Spatial-Temporal Routing Optimization for Last-Mile Food Delivery,'' 2025.

    \item Al-Mansouri F., Ibrahim H., Khalil A., ``Composite GiST Indexing for Spatiotemporal Query Optimization in High-Volume Logistics,'' 2026.

    \item Kavoosi A., Sadjadi S. J., Mohammadi M., ``Dynamic Pricing and Inventory Control of Perishable Products Using Deep Reinforcement Learning,'' 2026.

    \item Muller T., Schmidt K., Weber R., ``Blockchain-Based Food Traceability Framework with GDPR-Compliant Off-Chain Storage,'' 2025.

    \item Patil V., Joshi A., Deshmukh N., ``Smart Contract-Driven Food Donation Platforms with QR-Code Verification and Incentive Mechanisms,'' 2025.
\end{enumerate}
\end{frame}

%--------------------------------------------------------
% REFERENCES PAGE 4 (Papers 22--25)
%--------------------------------------------------------
\begin{frame}{References (4/4)}
\tiny
\begin{enumerate}
\setcounter{enumi}{21}
    \item Zhang H., Li M., Xu F., ``Explainable AI-Enhanced QR-Based Food Supply Chain Traceability for Consumer Trust,'' 2025.

    \item Sereda D., Kowalski M., ``Comparative Performance Analysis of Server-Sent Events vs WebSockets in High-Concurrency Systems,'' 2025.

    \item Rahman A., Hossain S., Karim T., ``Serverless Transaction Management for Event-Driven Food Delivery Microservices Using Saga Patterns,'' 2025.

    \item Park J., Kim S., Lee H., ``BE-YOLOv8: Improved YOLOv8 with Attention Mechanisms for Real-Time Food and Waste Classification,'' 2025.
\end{enumerate}
\end{frame}

%----------------------------------------------------------------------------------------
\end{document}
